\section{Related work}

We review existing fact-checking datasets. Then, we give an overview of related tasks for scientific document understanding. This work also provides a potential testbed for rationalized NLP models; see \citet{DeYoung2019ERASERAB} for an overview.

%%%%%%%%%%%%%%%%%%%%%%%%%%%%%%%%%%%%%%%%

\subsection{Fact checking} \label{sec:fact_checking}

\dw{Shorten} The proliferation of misinformation in politics and on the web has motivated the creation of a number of fact-checking datasets of varying size and annotation complexity. For a comprehensive overview, see \citet{Hanselowski2019ARA}. We discuss characteristics of fact-checking datasets that are relevant to our work by comparing \ours to three other large datasets annotated with sentence-level rationales: \fever \cite{Thorne2018FEVERAL}, \snopes \cite{Hanselowski2019ARA}, and Perspectrum \cite{Chen2019SeeingTF}.

\paragraph{Natural vs. synthetic.} As discussed in \citet{Hanselowski2019ARA}, \emph{natural} claims occur naturally on, for instance, political debate websites or Twitter. \snopes and Perspectrum have natural claims. \emph{Synthetic} claims are written during the construction of a fact-checking dataset by mutating a sentence from the evidence document. \fever claims are synthetic. Natural claims are more realistic, but more difficult to generate and verify. Claims in \ours are natural, since they are based on citations written by researchers about prior scientific work.

\paragraph{Factual vs. opinion-based.} The statements in \fever and \snopes are primarily factoid claims like ``\emph{Barak Obama was the $44^{th}$ President of the United States.}'' It is sensible to assign a single label indicating whether a corpus supports or refutes this claim. Claims in Perspectrum are opinions like ``\emph{Animals should have lawful rights}''. For these claims, it is inappropriate to assign a single global veracity label. Instead, individual documents agreeing or disagreeing with the claim should be identified.

Scientific claims have properties of both factoid and opinion-based claims. For example, ``\emph{The $R_0$ of the novel coronavirus is 2.5}'' makes a factual assertion, but its veracity may be the subject of active scientific research. Therefore, it is not be appropriate to assign a single truth label to this claim based on a corpus. Instead, like Perspectrum, we focus on identifying individual abstracts that support or refute a given claim.

\paragraph{Stance vs. entailment.}

\dw{Combine this with the factual vs. opinion-based section. Say we assume scientific articles do not deliberately misinform.}

The political fact-checking task is often divided into two sub-tasks \cite{derczynski-etal-2017-semeval}: \emph{stance detection} predicts the attitude of a document with respect to a claim which may be subjective and under dispute, while \emph{veracity prediction} evaluates the objective truth of the claim given that evidence documents may contain misinformation or be unreliable. For scientific fact-checking, we do not expect documents to contain deliberate misinformation, nor do we expect an abstract to express an opinionated stance. Like \fever, we focus on classifying the entailment relationship between a claim and a relevant abstract.


\begin{table}[t]
  \setlength{\tabcolsep}{.25em}
  \footnotesize
  \centering

    \begin{tabular}{ l  r  r  r }
      \toprule
      Dataset & \# Claims & Domain & Claim type \\
      \midrule
      \fever      & 185,445   & Wikipedia articles     & Synthetic             \\
      \snopes   & 6,422   & Political news      & Natural             \\
      \ours  & \numclaims  & Scientific text      & Natural             \\
      \bottomrule
    \end{tabular}

  \caption{Large fact-checking datasets annotated with rationale sentences. \dw{If keeping changes, add Perspectrum.}}
  \label{tbl:dataset_comparison}

\end{table}


%%%%%%%%%%%%%%%%%%%%

% Old way of doing things.

% The proliferation of misinformation in politics and on the web has motivated the creation of a number of fact-checking datasets of varying size and annotation complexity. For a comprehensive overview, see \citet{Hanselowski2019ARA}. We discuss the two largest fact-checking datasets for which sentence-level rationales are available. A summary is shown in Table \ref{tbl:dataset_comparison}. In our experiments, we compare the performance of models trained on these datasets to models trained on \ours.

% \fever \cite{Thorne2018FEVERAL} is a popular fact-checking dataset featuring factoid claims based on Wikipedia articles. \fever claims are \emph{synthetic}, meaning that they were created for a fact-checking dataset by altering sentences from the document containing the evidence. The claims in \ours are \emph{natural} -- they are based on citations written by researchers about prior scientific work.

% The claims in the \snopes dataset \cite{Hanselowski2019ARA} were collected from the Snopes political fact-checking website, and verified with sentence-level rationales from news articles. \snopes is more similar to \ours because its claims occur naturally in online political discussion. However, the scientific documents in \ours present additional linguistic challenges discussed in Section \todo.

% The political fact-checking task is often divided into two sub-tasks \cite{derczynski-etal-2017-semeval}: \emph{stance detection} predicts the attitude of a document with respect to a claim which may be subjective and under dispute, while \emph{veracity prediction} evaluates the objective truth of the claim given that evidence documents may contain misinformation or be unreliable. For scientific fact-checking, we do not expect documents to contain deliberate misinformation. Like \fever, we focus on \emph{recognizing textual entailment} of the claims with respect to evidence.

% \begin{table}[t]
  \setlength{\tabcolsep}{.25em}
  \footnotesize
  \centering

    \begin{tabular}{ l  r  r  r }
      \toprule
      Dataset & \# Claims & Domain & Claim type \\
      \midrule
      \fever      & 185,445   & Wikipedia articles     & Synthetic             \\
      \snopes   & 6,422   & Political news      & Natural             \\
      \ours  & \numclaims  & Scientific text      & Natural             \\
      \bottomrule
    \end{tabular}

  \caption{Large fact-checking datasets annotated with rationale sentences. \dw{If keeping changes, add Perspectrum.}}
  \label{tbl:dataset_comparison}

\end{table}





% The fact-checking task requires a system to verify a claim against a corpus. To verify factoid claims like ``\emph{Barak Obama was the $44^{th}$ president of the United States}'', a fact-checking system is required to output a single label indicating the claim's veracity given the corpus. The corpus \supports the claim when the claim is \emph{entailed} by a collection of sentences from the corpus. Most fact-checking datasets, including \fever and \snopes, take this approach. For opinion-based claims like ``\emph{Animals should have lawful rights}'', predicting a single label is not appropriate. Instead, the system should identify all relevant documents and label their \emph{stance} toward the claim. This is the approach used in the Perspectrum dataset \cite{Chen2019SeeingTF}.

%%%%%%%%%%%%%%%%%%%%%%%%%%%%%%%%%%%%%%%%

\subsection{Scientific document understanding} \label{sec:scidoc}

Although there are no previously existing datasets for scientific fact-checking, a number of scientific corpora exist on related tasks in scientific document understanding.

% \hanna{Also information extraction} \hanna{We do not need to add too much details on everything here, but it is important to say why fact checking is different from other problems? esp. QA}

\paragraph{Question answering} \dw{Much shorter} The BioASQ Task B \cite{Tsatsaronis2015AnOO} dataset consists of roughly 3,000 open-domain question / answer pairs issued by medical professionals, annotated with answers from PubMed articles. The questions come in a variety of formats, including extractive, ``yes / no'', and abstractive. PubMedQA \cite{Jin2019PubMedQAAD} is a large corpus of weakly-supervised ``yes / no'' questions over PubMed abstracts. The MEDIQA 2019 shared task \cite{Abacha2019OverviewOT} presents 40,000 questions from the related domains of clinical medicine and consumer health.

\paragraph{Citation contextualization}  The goal of \emph{citation contextualization} is to identify all spans in a cited document that are relevant to a particular citation in a citing document. A corpus of 20 biomedical documents with annotated contextualized citations was released at TAC 2014\footnote{\url{https://tac.nist.gov/2014/BiomedSumm/index.html}}. While the dataset was annotated by domain experts, the average inter-annotator agreement rate on annotated spans was only 21.7\% \cite{Cohan2015MatchingCT}.

\paragraph{Evidence inference}  Also related to fact-checking is the task of biomedical \emph{evidence inference} \cite{Lehman2019InferringWM}, which involves predicting the effect of a medical intervention on a specified outcome. Evidence inference can be viewed as a special case of claim verification.

\paragraph{Scientific information extraction} \dw{Shorter} Data and models exist to extract genes, proteins, and drugs \cite{Kim2003GENIAC, Krallinger2017OverviewOT}, computer science methods \cite{Luan2018MultiTaskIO}, and PICO clinical trial tags \cite{Nye2018ACW} from research abstracts. While our work focuses on verifying claims using article text, future efforts could incorporate information from automatically-extracted knowledgebases in the claim verification process.
